\documentclass{article}
\usepackage[utf8]{inputenc}
\usepackage[russian]{babel}
\usepackage{amsmath}
\usepackage{amssymb}
\usepackage{geometry}
\usepackage{pgfplots}
\geometry{a4paper, margin=1in}

\title{Лабораторная работа №4}
\author{}
\date{}

\begin{document}

\maketitle

\section{Контекстная свобода}

Для начала заметим, что атрибут по факту измеряет высоту возможного дерева разбора.
Пусть наперед задана строка с $n$ числом $a$, тогда атрибут лежит от $\lceil \log_2(n+1) \rceil$ до $n$.

Тогда условие эквивалентно $\lceil \log_2(T_1.attr+1) \rceil < T_2.attr$, а их атрибуты означают число букв $a$ в них буквально.
Отсюда получаем слово для накачки, которое пересекаем с соответствующим регулярным выражением.

Рассмотрим слово $(bba)^{2^n-1} bbaaabb (abb)^{n+1}$.

Понятно, что отдельно компоненты можно накачкой вывести из языка. Но также если и синхронно качать, то выйдем из языка тоже, ведь нулевой накачкой мы значение атрибута у левой части не изменим, так как логарифм берется вверх и такое вычитание не повлияет на атрибут, а правый уже уйдет и нарушится неравенство на атрибуты.

Поэтому язык не является контекстно свободным.


\section{Оптимизированный парсер}

Рассмотрим идею построения детерминированного алгоритма разбора.
Идем от листьев дерева разбора к его вершине.

Первоначально заметим, что подстрока \texttt{aaa} может быть получена только по правилу $S \rightarrow aaa$.
Поэтому на первом проходе сразу заменяем все вхождения \texttt{aaa} на $S$.
По аналогичным причинам сразу же заменяем все вхождения \texttt{SbS} на $S$ по правилу $S \rightarrow SbS$.

Далее работаем с третьим правилом $S \rightarrow bbSbb$.
Просматриваем строку слева направо.
Если встречаем подстроку \texttt{bbSbb}, то следует рассмотреть именно это правило.
В рамках этого правила будем использовать два указателя:
один будет считать возможное число соответствующих букв \texttt{a} справа от $S$,
другой~--- слева от $S$. Также из-за неравенства справа от $S$ хотя бы одна \texttt{a} должна быть, поэтому на правое вхождение такого же $S$ наслоения быть не должно.
При этом оба указателя движутся по максимально возможным последовательностям \texttt{b},
так как разные вхождения $S$ не могут пересекаться.
Если условие из описания языка (соотношение между атрибутами) выполняется,
то вся подстрока \texttt{bbSbb} заменяется на $S$;
если условие не выполняется, то данное вхождение пропускается.

После обработки очередного вхождения перемещаемся вправо, начиная именно с самого левого подходящего фрагмента.
Это вызвано принципом ``жадности'': левая часть неравенства на атрибуты должна быть как можно меньше,
чтобы не нарушить возможность разбора последующих частей строки.

Приоритет при замене отдается правилу $S \rightarrow SbS$;
только после того, как таких вхождений не осталось, переходим к обработке $S \rightarrow bbSbb$.
При этом, как уже сказано, каждый раз выбирается самое левое возможное вхождение.

В конечном счете, если исходная строка принадлежит языку, алгоритм приведет её к единичному символу $S$.

Таким способом строится детерминированный разбор,
который для каждой строки языка выбирает одно конкретное дерево разбора,
удовлетворяющее условиям на атрибуты, тем самым реализуя требуемую семантику языка.

\section{Оценка сложности оптимизированного парсера}

Эффективный парсер имеет кубическую сложность $O(n^3)$ в худшем случае.

Рассмотрим пооперационно:

\begin{enumerate}
    \item \texttt{replaceAll("aaa", "S")} — $O(n)$
    
    \item Внешний цикл \texttt{while} выполняется не более $O(n)$ раз, так как каждое успешное применение правила сокращает длину строки.
    
    \item \texttt{replaceAll("SbS", "S")} — $O(n^2)$ в худшем случае (проход по строке и создание новой строки)
    
    \item Внутренний цикл поиска паттернов для правила $S \rightarrow bbSbb$:
    \begin{itemize}
        \item Поиск подстрок \texttt{"bbSbb"} может выполняться $O(n)$ раз
        \item Каждый вызов функции проверки условия (функция \texttt{f}) — $O(n)$
    \end{itemize}
    Итого для этого этапа: $O(n) \times O(n) = O(n^2)$
\end{enumerate}

Общая оценка: 
\[
O(n) \times \big( O(n^2) + O(n^2) \big) = O(n^3)
\]

Таким образом, предложенный алгоритм имеет кубическую временную сложность в худшем случае.

\section{Тестирование}

Для проверки эффективности предложенного алгоритма было проведено тестирование на двух типах данных:
\begin{itemize}
    \item словах, принадлежащих языку;
    \item словах, не принадлежащих языку.
\end{itemize}

Для каждой длины измерялось время выполнения двух алгоритмов:
\begin{itemize}
    \item обычного парсера;
    \item оптимизированного парсера, описанного в разделе 2.
\end{itemize}

Результаты представлены на графиках ниже. (Значения времени приведены в миллисекундах.)

\begin{figure}[h!]
\centering
\begin{minipage}{0.48\textwidth}
\centering
\begin{tikzpicture}
\begin{axis}[
    title={Обычный парсер: слова из языка},
    xlabel={Длина строки $n$},
    ylabel={Время $t$, мс},
    xmin=0, xmax=100,
    ymin=0, ymax=151,
    xtick={30,42,59,75,99},
    ytick={0,10,50,100,150},
    legend style={at={(0.5,-0.2)}, anchor=north},
    grid=both,
    width=\textwidth,
    height=0.8\textwidth
]

% Обычный парсер для слов из языка (синий)
\addplot[
    color=blue,
    mark=triangle,
    mark size=3pt,
    thick
]
coordinates {
    (30, 1.305)(42, 9.339)(59, 10.256)(75, 9.727)(99, 150.88)
};
\addlegendentry{Обычный парсер}

\end{axis}
\end{tikzpicture}
\caption{Время разбора для слов из языка (обычный парсер)}
\label{fig:ordinary_in_lang}
\end{minipage}
\hfill
\begin{minipage}{0.48\textwidth}
\centering
\begin{tikzpicture}
\begin{axis}[
    title={Оптимизированный парсер: слова из языка},
    xlabel={Длина строки $n$},
    ylabel={Время $t$, мс},
    xmin=0, xmax=100,
    ymin=0, ymax=1,
    xtick={27,40,57,74,93},
    ytick={0,0.25,0.5,0.75,1},
    legend style={at={(0.5,-0.2)}, anchor=north},
    grid=both,
    width=\textwidth,
    height=0.8\textwidth
]

% Оптимизированный парсер для слов из языка (красный)
\addplot[
    color=red,
    mark=square,
    mark size=3pt,
    thick
]
coordinates {
    (27, 0.204)(40, 0.226)(57, 0.145)(74, 0.056)(93, 0.050)
};
\addlegendentry{Оптимизированный парсер}

\end{axis}
\end{tikzpicture}
\caption{Время разбора для слов из языка (оптимизированный парсер)}
\label{fig:optimized_in_lang}
\end{minipage}
\end{figure}

\begin{figure}[h!]
\centering
\begin{minipage}{0.48\textwidth}
\centering
\begin{tikzpicture}
\begin{axis}[
    title={Обычный парсер: слова не из языка},
    xlabel={Длина строки $n$},
    ylabel={Время $t$, мс},
    xmin=0, xmax=100,
    ymin=0, ymax=250,
    xtick={30,42,59,75,99},
    ytick={0,10,50,100,250},
    legend style={at={(0.5,-0.2)}, anchor=north},
    grid=both,
    width=\textwidth,
    height=0.8\textwidth
]

% Обычный парсер для слов не из языка (синий)
\addplot[
    color=blue,
    mark=triangle,
    mark size=3pt,
    thick
]
coordinates {
    (30, 3.322)(42, 8.093)(59, 20.947)(75, 54.827)(99, 211.695)
};
\addlegendentry{Обычный парсер}

\end{axis}
\end{tikzpicture}
\caption{Время разбора для слов не из языка (обычный парсер)}
\label{fig:ordinary_not_in_lang}
\end{minipage}
\hfill
\begin{minipage}{0.48\textwidth}
\centering
\begin{tikzpicture}
\begin{axis}[
    title={Оптимизированный парсер: слова не из языка},
    xlabel={Длина строки $n$},
    ylabel={Время $t$, мс},
    xmin=0, xmax=100,
    ymin=0, ymax=0.05,
    xtick={30,42,59,75,99},
    ytick={0,0.01,0.03,0.04,0.05},
    legend style={at={(0.5,-0.2)}, anchor=north},
    legend style={at={(0.5,-0.2)}, anchor=north},
    grid=both,
    width=\textwidth,
    height=0.8\textwidth
]

% Оптимизированный парсер для слов не из языка (красный)
\addplot[
    color=red,
    mark=square,
    mark size=3pt,
    thick
]
coordinates {
    (30, 0.025)(42, 0.017)(59, 0.028)(75, 0.035)(99, 0.031)
};
\addlegendentry{Оптимизированный парсер}

\end{axis}
\end{tikzpicture}
\caption{Время разбора для слов не из языка (оптимизированный парсер)}
\label{fig:optimized_not_in_lang}
\end{minipage}
\end{figure}

\end{document}